\begin{lemma}
\label{editorial-source-lemma-separably-closed-irreducible}
Let $k$ be a separably closed field.
Let $R$, $S$ be $k$-algebras. If $R$, $S$ have a unique
minimal prime, so does $R \otimes_k S$.
\end{lemma}

\begin{proof}
Choose the original perfect closure
$\overline k/k$ of Definition
\ref{editorial-source-definition-perfection}. It is
algebraically closed in this situation.
Indeed take an algebraic closure
$\Omega/k$. Its maximal separable
subextension is $k$, because $k$
is separably closed. Fields, Lemma
\ref{fields-lemma-separable-first}
therefore makes $\Omega/k$ purely
inseparable. The algebraically
closed field $\Omega$ is perfect,
so uniqueness of perfect closure
gives the unique $k$-isomorphism
$\overline k\simeq\Omega$.
In characteristic zero both
extensions are already the identity.

For each original $k$-algebra $A$,
the map $A\to A\otimes_k\overline k$
is injective by extension of a
vector-space basis. In positive
characteristic the receiving algebra
is generated by $1\otimes\lambda$,
each with a $p$-power in the image
of $k$, and the corresponding
integer multiple is zero.
Lemma \ref{editorial-source-lemma-p-ring-map} gives
a homeomorphism on spectra; in
characteristic zero the map is
the identity under multiplication.
Apply this to $R$, $S$ and the
original $R\otimes_k S$.
The scalar comparison is the
isomorphism
$$
(R\otimes_k S)\otimes_k\overline k
\longrightarrow
(R\otimes_k\overline k)\otimes_{\overline k}
(S\otimes_k\overline k),
\qquad
(r\otimes s)\otimes c\longmapsto
(r\otimes c)\otimes(s\otimes1).
$$
Its inverse sends
$(r\otimes a)\otimes(s\otimes b)$
to $(r\otimes s)\otimes ab$.
The original scalar actions prove
balancing and both composites.
The two changed factors still
have a unique minimal prime.
It therefore suffices to prove
the assertion over the actual
algebraically closed field
$\overline k$, and then descend
through these homeomorphisms.

Work now over an algebraically
closed field $k$, keeping the
algebras then under consideration.
For arbitrary $R,S$, the tensor
quotient
$$
R\otimes_k S\longrightarrow R_{red}\otimes_k S_{red}
$$
is surjective with kernel
the sum of the two ideals
generated by the original factor
nilradicals. The inverse of its
induced quotient isomorphism
sends $[r]\otimes[s]$ to the
class of $r\otimes s$.
Each element of this kernel is
a finite sum of nilpotent pure
tensors; if their exponents are
$d_1,\ldots,d_m$, the full
multinomial expansion at exponent
$1+\sum_i(d_i-1)$ is zero.
Thus the kernel is locally
nilpotent and Lemma
\ref{editorial-source-lemma-surjective-locally-nilpotent-kernel}
identifies the spectra.
Since the original factors have
unique minimal primes, their
reductions are nonzero domains.
We may consequently prove the
remaining assertion for these
two domains, without losing the
nilpotent information in the
original quotient map.

For domains $R,S$ over this
algebraically closed $k$, their
tensor ring is nonzero, as is
seen by extending the vectors
$1\in R$ and $1\in S$ to bases.
It is reduced by Lemma
\ref{editorial-source-lemma-perfect-reduced}.
If it were not a domain, there
would be nonzero original tensors
$z,w$ with $zw=0$. Each uses
finitely many coefficients from
$R$ and $S$. Choose the subalgebras
$R_0,S_0$ generated by all those
coefficients. They are domains
of finite type over the
algebraically closed field $k$.
The map $R_0\otimes_k S_0\to R\otimes_k S$
is injective by applying the
two vector-space injections
successively. Hence the same
tensors $z,w$ are nonzero already
in $R_0\otimes_k S_0$, and
their product is zero there.

The original $R_0$-algebra
$R_0\otimes_k S_0$ is flat:
a $k$-basis of $S_0$ gives a
free $R_0$-module basis.
It is finitely presented since
$S_0$ has a finite presentation
over $k$ by the Hilbert basis
theorem, and the same generators
and relations give its scalar
extension over $R_0$.
For each maximal ideal
$\mathfrak m\subset R_0$, the
Nullstellensatz (Theorem \ref{editorial-source-theorem-nullstellensatz}) gives the
original scalar isomorphism
$k\to R_0/\mathfrak m$.
The fibre map
$$
(R_0\otimes_k S_0)\otimes_{R_0}R_0/\mathfrak m
\longrightarrow S_0,
\qquad(r\otimes s)\otimes\overline a
\longmapsto\overline r\,\overline a\,s
$$
uses that isomorphism; its
inverse is $s\mapsto(1\otimes s)\otimes1$.
These fibres are irreducible.
The maximal ideals are dense
because $R_0$ is a finite type
algebra over a field and is
Jacobson, by Lemma
\ref{algebra-lemma-finite-type-field-Jacobson}.
Lemma \ref{editorial-source-lemma-flat-fibres-irreducible}
therefore makes the tensor's
spectrum irreducible. It is
also reduced by perfectness,
so its nilradical is the zero
prime and it is a domain.
This contradicts the original
nonzero $z,w$ with product zero.
Thus the arbitrary domain tensor
is a domain, and the preceding
quotient and field-extension
maps give the result for the
original $R,S$ over the original
separably closed field.
\end{proof}